\documentclass[11pt, a4paper]{article}

% Essential Packages
\usepackage[utf8]{inputenc}
\usepackage{geometry}
\geometry{margin=1in}
\usepackage{amsmath, amssymb, graphicx}
\usepackage{hyperref}

% We can solve parity game using algorithms that use BDDs, in earlier years we did this with fixpoint algorithms already. However these algorithms are quite vulnerable to particular parity games. Goal of the project is to study different symbolic versions of parity game algorithms, that are more robust, such as tangle learning, tangle attractors, fixpoints with justification, and possibly quasi-polynomial algorithms.

%It is relatively straight forward to start with justification fixpoints and see how far we can go. We have "fixpoints with freezing", and I think the next targets would be "fixpoints with justification" and "priority promotion".

\title{A symbolic approach to FPJ}
\author{Thomas van de Weerd \\ \small University of Twente}
\date{\today}

\begin{document}

\maketitle

\begin{abstract}
A concise summary of your research problem, proposed methodology, and expected impact (150–350 words).
\end{abstract}

\section{Introduction}
\label{sec:intro}
Describe the background, context, and the core problem you want to study. State why this topic matters right now.

\section{Literature Review \& Research Gap}
\label{sec:litreview}
Summarize previous work in the field and pinpoint what is missing (the research gap) that your project will address.

\section{Research Questions and Objectives}
\label{sec:objectives}
\begin{itemize}
    \item \textbf{Primary Question:} What is the main question you aim to answer?
    \item \textbf{Objective 1:} Specific goal one.
    \item \textbf{Objective 2:} Specific goal two.
\end{itemize}

\section{Proposed Methodology}
\label{sec:methodology}
Explain the tools, data collection, experiments, or theoretical frameworks you will use to conduct the research and why you chose them.

\section{Expected Timeline}
\label{sec:timeline}
Outline the major phases of the project (e.g., Month 1–3: Literature review and setup; Month 4–8: Data collection/experiments; Month 9–12: Analysis and writing).

\section{Significance and Impact}
\label{sec:impact}
Explain how your findings will contribute to academic knowledge or practical applications.

\begin{thebibliography}{9}
\bibitem{Author2026}
Author, A. (2026). \emph{Sample Journal Article Title}. Journal Name, 12(3), 45-60.
\end{thebibliography}

\end{document}
